\documentclass[11pt]{article}
\usepackage[margin=2.6cm]{geometry}
\usepackage{amsmath,amssymb,amsthm,booktabs,hyperref}
\newtheorem{proposition}{Proposition}
\newtheorem{theorem}{Theorem}
\newtheorem{conjecture}{Conjecture}
\title{Dilated Singer sets and the excess in Erd\H{o}s's Sidon problem:\\ a curve-counting theorem, forbidden-structure laws,\\ a resonance theorem, and a reduction}
\author{Vichien Fugsukjit\\ Independent researcher, Bangkok\\ \small(designed and carried out with AI assistance: Anthropic Claude)}
\date{Version 3, 21 October 2026 (Version 2: 19 October 2026; Version 1: 16 October 2026)}
\begin{document}
\maketitle
\begin{abstract}
Let $h(N)$ be the maximal size of a Sidon set in $\{1,\dots,N\}$. Erd\H{o}s asked whether $h(N)=N^{1/2}+O(1)$ (Guy C9); the weaker
$h(N)=N^{1/2}+O_\varepsilon(N^\varepsilon)$ is Erd\H{o}s problem \#30. We study random dilates of Singer perfect difference sets $S$
($|S|=q+1$, modulus $m=q^2+q+1$) and the number $Z$ of points of a dilate in a window of length $Cq$.
\emph{Proved:} (i) if the factorial moments of $Z$ converge to the Poisson ones for every fixed order and every window constant $C$, then $h(N)-N^{1/2}$ is unbounded
(Proposition~\ref{prop:red}); (ii) for a ratio $\rho=u/v$, the number $N(\rho)$ of triples of $S$ with $c-a=\rho(b-a)$ is the number
of $\mathbb F_q$-points of a twisted form of a curve $C_{u,v}$ that does not depend on $q$, so $N(\rho)/E-1=O_\rho(q^{-1/2})$ whenever
$C_{u,v}$ minus its diagonal is absolutely irreducible, which we certify by exact computation for 17 ratio orbits, including all
non-forbidden ratios of height $\le4$ (Theorem~\ref{thm:curve}); (iii) the identity $\sum_\rho N(\rho)^2=\sum_{\lambda,\mu}(|S\cap(\lambda S+\mu)|)_3$.
\emph{Observed, with pre-registered blind tests:} the deviation of the $k$-th factorial moment from the random value is, to first order,
$-[\binom k3P_3+\binom k4P_4]$, where $P_3,P_4$ are the window probabilities of 3-term progressions and parallelograms (asymptotically
$-[\tfrac32\binom k3+2\binom k4]/(Cq)$), for $k\le8$ (up to $k=7$ at $C=2$, $k=8$ at $C=4$) on 129 primes not previously used for $k\ge4$; and $\sum_\rho(N(\rho)-E)^2$ equals the order of the multiplier
group times $\sum_\rho E$, to within $0.8\%$ for each of 37 primes $293\le q\le2969$ with $m$ prime (mean $2.998$),
and within $10\%$ for 21 of 23 prime powers $121\le q\le2197$. \emph{New in Version 3:} $3k\mid N(\rho)$ (a free multiplier action, partly machine-checked in Lean), and a resonance theorem for cube $q=r^3$. For the six ratios of the $S_3$-orbit of $r+r^{-1}$ (roots modulo $m$ of $x^3-3x+1$ and $x^3-3x^2+1$), every point of the Singer line not fixed by the $q$-Frobenius has exactly $r$ partners, so $N(\rho)=r(q+1-e)$ with $e\in\{0,1\}$ ($m$ prime). This explains most of the excess of the variance ratio for cube $q$. There is one pre-registered blind test of the predicted excess ($q=12167$), and pre-registered tests of the resonant orbit and of $M/\mu\approx r$ ($q=4913,6859,15625$). The remaining step towards Erd\H{o}s's question is a uniform estimate
(an average over dilates), which we isolate but do not prove; the Erd\H{o}s question remains open.
\end{abstract}

\section{Background}
A set $A\subset\mathbb Z$ is Sidon if all sums $a+b$ ($a\le b$) are distinct. Erd\H{o}s and Tur\'an \cite{ET} showed
$h(N)\le N^{1/2}+O(N^{1/4})$; Singer's construction \cite{Si} with the Baker--Harman--Pintz prime-gap theorem \cite{BHP} gives
$h(N)\ge N^{1/2}-O(N^{0.2625})$; the constant in the $N^{1/4}$ term has since been improved several times (Balogh--F\"uredi--Roy
\cite{BFR} and later work listed at \cite{EP}). Erd\H{o}s offered USD~500 for deciding whether
$h(N)=N^{1/2}+O(1)$ \cite{Guy,MO}. The site \cite{EP} lists the weaker question $h(N)=N^{1/2}+O_\varepsilon(N^\varepsilon)$ as
problem \#30 (USD~1000), with the $O(1)$ form as a remark; the evidence below concerns the $O(1)$ question only.
A Sidon set $A$ modulo $m$ with $|A|=k$, cut at its largest cyclic gap $G$, gives a Sidon set of integers in an interval of length
$m-G$; for Bose sets \cite{Bo} ($m=q^2-1$, $k=q$) this gives $k-\sqrt{N}\approx G/(2q)$. Martin and O'Bryant \cite{MO} note that
the question would be answered (excess unbounded) if Bose sets had gaps that are not $O(q)$. Dilating $A$ by a unit $t$ preserves
the Sidon property and moves the gaps; the use of dilates was suggested in B.~Green's comment on \cite{Go} (see also \cite{Gr}; earlier computer searches of dilates and cuts: \cite{Z,RD}), and Shakan \cite{Sh}
proved that every $A\subset\mathbb F_p$ with $|A|>1$ has a dilate with a gap of at least $2p/|A|-2$.

\section{Dilate gaps of Bose sets}
For every prime $q$ in the ranges below we built $\mathrm{Bose}(q)\subset\mathbb Z/(q^2-1)$, checked the Sidon property, and computed
$G^*(q)=\max_t G(tA)$ over one multiplier per class $\{\pm t,\pm tq\}$. Positive control: $q=829$ gives excess $10.2797$, matching
Zhang's bound $K(p_{145})>10.279$ ($p_{145}=829$) \cite{Z}.
\begin{itemize}
\item $101\le q\le2003$ (279 primes): a $\log q$ law fitted on $q\le1000$, $g^*\approx-6.92+3.63\ln q$, gives held-out ($q>1000$)
RMSE $1.22$, against $4.25$ for a constant fitted on the same data. (Saturating bounded models were not tested.)
\item Blind test (law registered before computing): 48 primes $2011\le q\le5881$ (every 10th prime in $(2003,6000]$): RMSE $1.52$
(log) against $7.34$ (constant), mean signed error $-0.57$. The largest $q$ ($5881$) gives a Sidon set of size $5881$ in $[1,34{,}453{,}078]$ with excess $11.33$; the largest excess in
this set is $13.18$ ($q=5039$).
\item Null model: random $q$-subsets of $\mathbb Z/(q^2-1)$ show the same growth (Bose slightly lower at 15 of 20 primes), so the
$\log q$ growth is random-set behaviour; counts of 3-term relations $e_1L_1+e_2L_2+e_3L_3\equiv0$ among the discrete logarithms are at
the random level (mean per-$q$ ratio $0.983$ over 11 primes $101\le q\le383$; percentile bootstrap over the primes, 95\% interval
$[0.96,1.00]$).
\end{itemize}

\section{The moment route for Singer sets}
Let $S$ be a Singer set \cite{Si} ($|S|=n=q+1$, $m=q^2+q+1$), $\ell=Cq$, and $Z_{t,s}=\#\{a\in S:(ta+s)\bmod m\in[0,\ell)\}$ for
units $t$ and shifts $s$. Let $F_k$ be the average of the falling factorial $(Z)_k$ and $H_k=(n)_k(\ell)_k/(m)_k$ its value for a
uniformly random $n$-subset; $R_k=F_k/H_k-1$.
\begin{proposition}
For every unit $t$, $\operatorname{Var}_s Z_{t,s}=q\ell(m-\ell)/m^2$ and $R_2=0$. If $m$ is prime, averaging over all units,
\[ R_3=\sum_{\rho\ne0,1}\Bigl(\frac{N(\rho)}{E}-1\Bigr)\frac{A(\rho)}{\sum_\sigma A(\sigma)},\qquad
N(\rho)=\#\{(a,b,c)\in S^3\ \text{distinct}: c-a\equiv\rho(b-a)\},\ E=\frac{(n)_3}{m-2},\]
where $A(\rho)$ is the average over $u\in(\mathbb Z/m)^\times$ of the number of windows $[s,s+\ell)$ containing $\{0,u,u\rho\}$, and
$\sum_\sigma A(\sigma)=\ell(\ell-1)(\ell-2)/(m-1)$.
\end{proposition}
(The first part holds because every nonzero residue is exactly one difference of $tS$ and $(n)_2=m-1$; for the second, $m$ prime
lets every triple be written as $(a,b,c)$ with $c-a=\rho(b-a)$, and averaging over all units $u$ sends it to every pair $(u,u\rho)$.) Since a perfect difference set
has $N(-1)=N(2)=N(1/2)=0$ and $A(\rho)\approx\ell^2/(2m)$ for these $\rho$, the three forbidden ratios contribute
$-3/(2Cq)(1+O(1/q))$ (within $0.2\%$ for $q\ge293$, $1.2\%$ at $q=41$).
\begin{itemize}
\item $R_k\sim c_k(C)/q$ (moments estimated from 4000 sampled multipliers): registered blind test for $k=4$ on 11 new primes
$409\le q\le601$ ($C=2$: RMSE $0.0028$ for $c/q$ against $0.0238$ for a constant; $C=4$: $0.0013$ against $0.0112$); the same is seen
descriptively for $k=3,5,6$, with $|c_k|$ growing quickly in $k$.
\item Exact over all units: registered tests on 18 primes $293\le q\le1373$ ($m$ prime) and, from a sealed data set opened under a
frozen plan, 17 primes $1559\le q\le2969$: the mean of $C\,qR_3$ is between $-1.45$ and $-1.41$, resp.\ between $-1.47$ and $-1.46$,
inside the registered bands around $-3/2$ for all $C\in\{2,3,4,5,8\}$. The first test could not separate $-3/2$ from an earlier
empirical value $-1.17$; the second band ($[-1.59,-1.31]$) excludes it.
\item The RMS deviation $D_{\rm mid}$ of $N(\rho)/E$ over ratios of height $3$--$6$: the registered decay test on $293\le q\le1373$
alone \emph{failed} (slope $-0.13\pm0.17$, outside the band $[-0.65,-0.35]$; few independent values per $q$). A second test, frozen
before opening the sealed data but pooling them with the already seen first set, gives slope $-0.504\pm0.098$ over $293\le q\le2969$
(the sealed set alone: $-1.28\pm0.39$, i.e.\ noisy); $\sqrt q\,D_{\rm mid}$ lies between $0.49$ and $1.84$. Over all ratios the exponent is $-0.50$ in both
sets, with $\sqrt q\,D_{\rm all}\approx1.73$ for Singer sets against about $1.0$ for random sets, consistent with counting noise
inflated by the size-3 orbits of the multiplier $q$. These data are consistent with, but do not establish, a $q^{-1/2}$ law.
\end{itemize}
Sections~4--8 and~\ref{sec:rem} (new in Version~2; Sections~\ref{sec:orb}--\ref{sec:resblind} are new in Version~3) turn the heuristic of Version~1 into a theorem for each fixed ratio, give exact identities,
explain the constants $c_k$, and state the reduction precisely.

\section{The reduction}\label{sec:red}
\begin{proposition}\label{prop:red}
Fix $C>0$ and let $\ell=\lfloor Cq\rfloor$. Suppose that for every fixed $k\ge1$, $R_k=R_k(q,C)\to0$ as $q\to\infty$ through primes.
Then for all large $q$ some dilate of $S$ has an empty arc of length $\ell$, and $\limsup_{N\to\infty}(h(N)-N^{1/2})\ge(C+1)/2$.
In particular, if this holds for every $C$, then $h(N)-N^{1/2}$ is unbounded.
\end{proposition}
\begin{proof}
Choose $(t,s)$ uniformly. Since $Z\le n$, inclusion--exclusion gives $\Pr(Z=0)=\sum_{k\ge0}(-1)^kF_k/k!$, and the Bonferroni
inequalities give $\Pr(Z=0)\ge\sum_{k=0}^{2K+1}(-1)^kF_k/k!$ for every $K$. Since $H_k=(n)_k(\ell)_k/(m)_k\to C^k$, the hypothesis gives
$F_k\to C^k$ for each fixed $k$. Choose $K$ with $\bigl|\sum_{k\le2K+1}(-1)^kC^k/k!-e^{-C}\bigr|<e^{-C}/2$; then $\Pr(Z=0)>0$ for all
large $q$, so some dilate $tS$ misses a window of length $\ell$. Rotating, $tS$ lies in an interval of length $m-\ell$; a Sidon set
modulo $m$ is a Sidon set of integers, so we get $q+1$ elements in $[1,N]$ with $N=m-\ell=q^2+(1-C)q+O(1)$ and
$q+1-N^{1/2}=(C+1)/2+O(1/q)$.
\end{proof}
The cutting step is classical (Section~1); what we add is the reduction of unbounded excess to the convergence of fixed-order factorial
moments of window counts. We did not find it in the literature (Section~\ref{sec:lim}).

\section{Universal curves: a theorem for each fixed ratio}
Write $S=\{i:g^i\in L\}$ with $g$ a generator of $\mathbb F_{q^3}^\times$ and $L$ a 2-dimensional $\mathbb F_q$-subspace. For
$\rho=u/v$ put $(\alpha,\beta,\gamma)=(u-v,v,-u)$, so $\alpha+\beta+\gamma=0$, and let
\[C_{u,v}:\qquad s^\alpha r^\beta t^\gamma=1,\qquad (s-1)^\alpha(r-1)^\beta(t-1)^\gamma=1 .\]
\begin{theorem}\label{thm:curve}
(a) For $q$ prime, $N(\rho)$ equals the number of $\mathbb F_q$-points of a twisted form of $C_{u,v}$ off its diagonal $s=r=t$, up to
$O_\rho(1)$; the twist is Frobenius composed with the order-3 map $s\mapsto1-1/s$.
(b) For $\rho\in\{-1,2,1/2\}$, $C_{u,v}$ is the diagonal, so $N(\rho)=0$.
(c) If $C_{u,v}$ minus the diagonal is absolutely irreducible over $\overline{\mathbb Q}$, then $N(\rho)/E-1=O_\rho(q^{-1/2})$ for all
primes $q>q_0(\rho)$.
(d) The hypothesis of (c) holds for every ratio of height $\le4$ other than $-1,2,1/2$ (five orbits) and for the orbits (under $\rho\mapsto1-\rho,1/\rho$) of
$-5,-5/2,-5/4,-6,-6/5,-7,-7/2,-7/6,-8,-8/3,-8/5,-8/7$: 17 orbits in all. With (b), 18 of the 22 orbits of height $\le8$ are settled.
\end{theorem}
\begin{proof}[Proof sketch]
(a) With $x=g^a$, $y=g^b$, $z=g^c$ the condition is $x^{u-v}z^{v}y^{-u}\in\mathbb F_q^\times$. Parametrise $P(L)\cong\mathbb P^1(\mathbb F_q)$
by $x=e+sf$; the three conjugate linear forms $e^{q^i}+sf^{q^i}$ have distinct roots, and a M\"obius change of variable over
$\mathbb F_{q^3}$ sending them to $\infty,0,1$ turns $w\in\mathbb F_q^\times$, i.e.\ $w=w^q=w^{q^2}$, into the two equations of
$C_{u,v}$ (the constants cancel because $\alpha+\beta+\gamma=0$); Frobenius becomes the stated twist. (b) For exponents $(-2,1,1)$ the
equations give $r+t=2s$ and $rt=s^2$. (c) Absolute irreducibility over $\overline{\mathbb Q}$ persists modulo all large primes; apply the
Weil bound for singular curves \cite{AP} to a $p$-independent projective model; excluded points are $O_\rho(1)$ and $E=q+O(1)$.
(d) Eliminating one variable by a resultant gives a plane curve $G$; we exhibit smooth points of $G$ over number fields whose degrees
have $\gcd1$ (or one quadratic or cubic field over which $G$ stays irreducible). A geometric component through a smooth $K$-point is
defined over a subfield of $K$, so the number of components divides every such degree and equals $1$. Birationality of the projection
is checked by subresultants. Exact arithmetic in SymPy; certificates and code are released. An independent review checked (a)--(c)
and the certificate logic, and confirmed (a) by exact counts for $q=5,7,11,13$.
\end{proof}
Numerically, the genus of $C_{u,v}$ appears to grow at least quadratically in the height (for example $(n-3)(n-4)/2$ for exponents $(-n,1,n-1)$, $n\le6$,
by Riemann--Hurwitz with numerical monodromy), so Theorem~\ref{thm:curve} alone does not give $R_3\to0$; see Section~\ref{sec:rem}.
Four further orbits ($-5/3,-7/3,-7/4,-7/5$) exceeded our 10-hour certification budget; none failed.

\section{Exact identities}
Let $I_\lambda(\mu)=|S\cap(\lambda S+\mu)|$ and let $m$ be prime. Two ordered triples of $S$ with the same ratio determine a unique affine map
$x\mapsto\lambda x+\mu$ taking one to the other, hence
\[\sum_{\rho\ne0,1}N(\rho)^2=\sum_{\lambda\in(\mathbb Z/m)^\times}\ \sum_{\mu}\bigl(I_\lambda(\mu)\bigr)_3,\]
which we also checked exactly for $q=5,17,41,59$. Since every nonzero residue is a unique difference, $\sum_\mu(I_\lambda(\mu))_2=m-1$
for every $\lambda$: second moments are exactly Poisson. For the multipliers $\lambda\in\{1,q,q^2\}$, $\lambda S$ is a translate of
$S$ and one $\mu$ gives $I=n$, contributing exactly $3(n)_3$. Under $x\mapsto x^\lambda$, a dilate by $\lambda$ whose base-$q$ digits
have sum $d$ maps the line $L$ into a curve of degree $d$, so $I_\lambda\le d$ by B\'ezout unless $L$ lies on that curve (this does not happen for $d=2$ and odd $q$); for $d=2$ (for example $\lambda=-1$, a
conic) $(I_\lambda)_3=0$. Numerically, the non-multiplier $\lambda$ with $\sum_\mu(I_\lambda(\mu))_3=0$ are exactly the nine elements
$\{-1,2,1/2\}\cdot\{1,q,q^2\}$ for $q=41,59,167,173,293,383$ (all $\lambda$ computed); the arc and oval facts behind $-1$ and $2$ are classical
\cite{JV,Pi}.

\section{Higher moments: the forbidden-structure law (blind tests)}
A Sidon set contains no 3-term progression and no non-trivial parallelogram $x+y=z+w$. Let $P_3(\ell)$, resp.\ $P_4(\ell)$, be the
probability that an ordered distinct triple of positions in $[0,\ell)$ contains a 3-term progression, resp.\ that an ordered distinct
quadruple is a parallelogram; $P_3\sim3/(2\ell)$ and $P_4\sim2/\ell$. The first-order prediction is
\[R_k\approx-\Bigl[\tbinom k3P_3(\ell)+\tbinom k4P_4(\ell)\Bigr]\sim-\frac{\tfrac32\binom k3+2\binom k4}{Cq}.\]
Two registered blind tests (blind for $k\ge4$ in the first, for all $k$ in the second) used plans frozen with a recorded SHA-256 before computing, 3000 sampled unit classes per $q$, and the
inverse-variance pooled ratio measured/predicted. On 65 primes $607\le q\le1499$ with $C=2$, $k=3,\dots,6$ gave $1.02,1.06,1.06,1.05$
(all 95\% intervals contain $1$). On the other 64 primes of that range, $C=2$ gave $0.97,0.97,0.97,0.97,1.01$ for $k=3,\dots,7$ (the
$k=8$ cell is void because its null control failed), and $C=4$ gave $0.95,0.98,1.00,1.00,1.00,1.01$ for $k=3,\dots,8$. Random sets were
within $3$ standard errors of $0$ in every other cell. So no new forbidden structure appears up to $8$ points, and $c_k\sim k^4$ explains the fast growth
noted in Version~1.

\section{Variance law and prime powers}
\begin{conjecture}\label{conj:var}
$\sum_{\rho}(N(\rho)-E)^2=|G|\,\sum_\rho E\,(1+o(1))$, where $|G|$ is the order of the multiplier group ($|G|=3$ for prime $q$, $|G|=3k$ for
$q=p^k$), the sums run over $\rho$ with $\rho$ and $\rho-1$ units, and $E$ is the mean of $N(\rho)$ over these $\rho$
(equal to $(n)_3/(m-2)$ when $m$ is prime).
\end{conjecture}
By Proposition~\ref{prop:orb} (Section~\ref{sec:orb}), $3k\mid N(\rho)$ for every prime power $q=p^k$. For prime $q$ with $m$ prime, the ratio has mean $2.998$ (s.d.\ $0.008$, range $[2.979,3.019]$) over 37 primes
$293\le q\le2969$. For prime powers (Singer sets built in $\mathbb F_{p^{3k}}$; registered test) the ratio divided by $3k$ is given in
Table~\ref{tab:pp}. The per-$\lambda$ errors $\sum_\mu(I_\lambda(\mu))_3$ minus the random value have mean about $0$; excluding the nine
null $\lambda$ they are close to Gaussian, with standard deviation about $4.2\,q$ ($q=383$: $4.21q$, kurtosis $3.01$).
\begin{table}[h]\centering\small
\begin{tabular}{rrrrrr}\toprule $q$&$p^k$&$3k$&good $\rho$&ratio&ratio$/3k$\\ \midrule
121&$11^{2}$&6&2975&6.13&1.022\\
125&$5^{3}$&9&14059&9.91&1.101\\
128&$2^{7}$&21&11725&20.30&0.967\\
169&$13^{2}$&6&9145&6.53&1.089\\
243&$3^{5}$&15&50149&14.62&0.974\\
256&$2^{8}$&24&13145&24.35&1.014\\
289&$17^{2}$&6&16775&6.32&1.053\\
343&$7^{3}$&9&37135&10.62&1.180\\
361&$19^{2}$&6&30625&6.12&1.020\\
512&$2^{9}$&27&262655&28.04&1.039\\
529&$23^{2}$&6&55055&5.75&0.959\\
625&$5^{4}$&12&86855&12.05&1.004\\
729&$3^{6}$&18&449225&18.76&1.042\\
841&$29^{2}$&6&192335&6.00&1.000\\
961&$31^{2}$&6&195755&6.00&0.999\\
1024&$2^{10}$&30&245105&29.96&0.999\\
1331&$11^{3}$&9&1772891&9.48&1.053\\
1369&$37^{2}$&6&386425&6.07&1.012\\
1681&$41^{2}$&6&937945&6.00&1.001\\
1849&$43^{2}$&6&947903&5.98&0.996\\
2048&$2^{11}$&33&2997385&33.22&1.007\\
2187&$3^{7}$&21&4048957&21.03&1.001\\
2197&$13^{3}$&9&1609667&9.79&1.088\\
\bottomrule\end{tabular}
\caption{Variance ratio $\sum(N-E)^2/\sum E$ over $\rho$ with $\rho,\rho-1$ units, for Singer sets with $q=p^k$, $k\ge2$ (registered test: within 10\% of $3k$ for at least 80\% of $q\ge121$; here 21 of 23; the two misses and most of the excess at cube $q$ are accounted for by the resonance of Section~\ref{sec:res} (Table~\ref{tab:res}); for $q<121$, outside the registered range, the ratio$/3k$ is far from 1, e.g.\ $0.20$ at $q=9$, $0.69$ at $q=49$, $0.83$ at $q=81$; later runs: all 11 later non-cube $q$ ($2209,2401,2809,3125,3481,3721,4489,6561,8192,14641,16384$) give ratio$/3k\in[0.9980,1.0026]$; the later cube $q=4096,4913,6859$ give $1.022,1.038,1.064$, see Table~\ref{tab:res}).}\label{tab:pp}\end{table}
\begin{proposition}[conditional]
If $\sum_\rho(N(\rho)-E)^2\ll q^3$, then $R_3=O\bigl(\sqrt{\log q/q}\,\bigr)$.
\end{proposition}
\begin{proof}[Proof sketch]
By the Proposition of Section~3, $R_3=\sum_\rho(N/E-1)(w(\rho)-\bar w)$ with $w=A/\sum A$, because $\sum_\rho N(\rho)=(n)_3$ makes the
uniform part vanish. Cauchy--Schwarz gives $|R_3|\le E^{-1}\bigl(\sum_\rho(N-E)^2\bigr)^{1/2}\|w-\bar w\|_2$, and lattice-point counting
for $\{(x,y):y\equiv\rho x\}$ in $(-\ell,\ell)^2$ gives $\|w-\bar w\|_2\ll\sqrt{\log q}/q$.
\end{proof}

\section{Orbits of the multiplier group}\label{sec:orb}
\begin{proposition}\label{prop:orb}
Let $q=p^k$ and let $N(\rho)$ count triples $(a,b,c)\in S^3$ with $b-a$ a unit and $c-a=\rho(b-a)$ (for prime $m$ this is the count of Section~3).
Then $3k\mid N(\rho)$ for every $\rho$.
\end{proposition}
\begin{proof}
By the multiplier theorem, some translate of $S$ is fixed by multiplication by $p$ \cite{Ha,MR}. $N$ is translation invariant, so assume $pS=S$.
We have $p^{3k}-1=(q-1)m$. If $p^j\equiv1\pmod m$ with $0<j<3k$, then $m\mid p^d-1$ for $d=\gcd(j,3k)\le 2k$, which is impossible because $0<p^d-1<m$.
So $p$ has order $3k$ modulo $m$. The map $(a,b,c)\mapsto(pa,pb,pc)$ preserves the set of triples. A triple fixed by $p^j$, $0<j<3k$, would give
$(p^j-1)(b-a)=0$ with $b-a$ a unit, which is impossible. So the action is free and every orbit has $3k$ elements.
\end{proof}
This is elementary and probably known. The order statement and the freeness step are machine-checked in Lean~4 with Mathlib
(files \texttt{C9Order.lean}, \texttt{C9Orbit.lean}; standard axioms only). The multiplier theorem itself is not formalised.
So Conjecture~\ref{conj:var} says that $M=N/3k$ is Poisson-like. This holds for non-cube $q$ (Table~\ref{tab:pp}) but fails for cube $q$; the next two sections explain most of this failure.

\section{A resonance theorem for cube $q$}\label{sec:res}
Let $q=r^3$, where $r$ is a prime power. Let $F=\mathbb F_q$ and $K=\mathbb F_{q^3}=\mathbb F_{r^9}$. Write $\tau(u)=u^r$ and $\varphi=\tau^3$ (the $q$-Frobenius).
Let $L=\ker\mathrm{Tr}_{K/F}$, so that $S=\{i: g^i\in L\}$. Since $\tau$ commutes with $\varphi$, we have $\tau(L)=L$. Moreover $\varphi$ is $F$-linear on $L$ with $\varphi^3=1$.
Since $r^9=q^3\equiv1$ and $r^3=q\not\equiv1\pmod m$, $r$ has order $9$ modulo $m$. Put $\rho_1=r+r^{-1}\bmod m$.
\begin{theorem}\label{thm:res}
Let $x\in P(L)$ be a point of the Singer line.
\begin{itemize}
\item If $\varphi x\ne x$ (projectively), then there are exactly $r$ points $y\in P(L)\setminus\{x\}$ with $z/x=(y/x)^{\rho_1}$ for some $z\in P(L)$.
\item If $\varphi x=x$, there are none.
\end{itemize}
The number $e$ of $\varphi$-fixed points of $P(L)$ is $0$, $1$ or $2$ according as $r\equiv2$, $0$ or $1\pmod3$.
Consequently, if $m$ is prime, $N(\rho)=r\,(q+1-e)$ for each of the six ratios $\rho$ in the $S_3$-orbit of $\rho_1$; this orbit contains $\rho_0=r^4+r^{-4}=(r+1)r^4$.
If $r\equiv1\pmod3$ then $3\mid m$, so for prime $m$ we have $e\in\{0,1\}$, and $e=1$ exactly when $p=3$.
\end{theorem}
\begin{proof}
(i) \emph{Reduction.} Put $u=y/x$. Then $y\in P(L)$ iff $\mathrm{Tr}(xu)=0$. If $x$ is not $\varphi$-fixed, then $\psi=\varphi(x)/x\notin F$ and
$\mathrm{Tr}(x\psi)=\mathrm{Tr}(\varphi x)=0$, so $\{u:\mathrm{Tr}(xu)=0\}=F+F\psi$.
In $K^\times/F^\times$ the exponent $\rho_1$ acts as $u\mapsto\tau(u)\tau^{-1}(u)$, so we need $z=x\,\tau(u)\,\tau^{-1}(u)\in F^\times L$. That is, $\mathrm{Tr}(x\,\tau(u)\,\tau^{-1}(u))=0$.
Take $u=1+s\psi$ with $s\in F$; on $F$ we have $\tau s=s^r$ and $\tau^{-1}s=s^{r^2}$. Using $\mathrm{Tr}(x)=0$, the condition reads $As^r+Bs^{r^2}+Cs^{r+r^2}=0$, where
$A=\mathrm{Tr}(x\tau\psi)$, $B=\mathrm{Tr}(x\tau^{-1}\psi)$ and $C=\mathrm{Tr}(x\,\tau\psi\,\tau^{-1}\psi)$ lie in $F$.
Remove $s=0$ (that is, $y=x$) and put $Z=s^{-r}$, with $s=\infty$ corresponding to $Z=0$. This gives the affine $\mathbb F_r$-linear equation
\[A Z^r+BZ+C=0,\qquad Z\in F.\]
(ii) \emph{Two identities.} Write $\xi_j=\tau^j(x)$ for $j\bmod 9$, with $\xi_j+\xi_{j+3}+\xi_{j+6}=0$. Then $A$, $B$, $C$ are rational functions of $\xi_0,\dots,\xi_5$
with integer coefficients and monomial denominators, and $\tau$ acts as the index shift. Let $D$ be the Dickson matrix of $f(Z)=AZ^r+BZ$. Exact computation in Sage gives
\[N_{F/\mathbb F_r}(A)+N_{F/\mathbb F_r}(B)=0\quad\text{and}\quad \mathrm{adj}(D)\,(C,C^r,C^{r^2})^{T}=0.\]
After clearing the (nonzero) denominators, both are identities between polynomials with integer coefficients that hold over $\mathbb Q$. They therefore hold in every characteristic.
They were also re-checked directly over $\mathbb F_2,\mathbb F_3,\mathbb F_5,\mathbb F_7$.
(iii) \emph{Non-degeneracy.} The numerator of $A$ factors as $(\xi_1\xi_3-\xi_0\xi_4)(\xi_1^2+\xi_1\xi_4+\xi_4^2)$, and the denominator is $-\xi_1\xi_4\xi_7\ne0$.
The second factor vanishes iff $\psi$ is a root of $t^2+t+1$. The first vanishes iff $\tau x=cx$ with $c\in F$, which gives $\varphi x=N_{F/\mathbb F_r}(c)\,x$.
In both cases $x$ is $\varphi$-fixed. So $A\ne0$, hence $N(B)=-N(A)\ne0$ and $\mathrm{rank}\,D=2$.
(iv) \emph{Counting.} $N(-B/A)=N(-1)\,N(B)/N(A)=(-1)^{1+r+r^2}\cdot(-1)=1$, so the kernel of $f$ has exactly $r$ elements and $|\mathrm{Im}\,f|=r^2$.
Since $\mathrm{adj}(D)\,D=\det(D)\,I=0$, each row $(w_0,w_1,w_2)$ of $\mathrm{adj}(D)$ annihilates $(b,b^r,b^{r^2})$ for every $b\in\mathrm{Im}\,f$.
A nonzero row gives a nonzero linearised polynomial of degree at most $r^2$. Its zero set therefore has at most $r^2$ elements and contains $\mathrm{Im}\,f$, so it equals $\mathrm{Im}\,f$.
By (ii), $-C\in\mathrm{Im}\,f$, so there are exactly $r$ solutions $Z$, that is, $r$ partners $y$.
(v) \emph{Fixed points.} The $\varphi$-fixed points of $P(L)$ are the eigenlines of the $F$-linear map $\varphi|_L$, whose eigenvalues are cube roots of unity in $F$.
The eigenvalue $1$ needs $x\in F$, and then $\mathrm{Tr}(x)=3x$ vanishes only if $p=3$.
\begin{itemize}
\item If $r\equiv2\pmod3$, then $F$ has no primitive cube root of unity, so $e=0$.
\item If $p=3$, then $\varphi|_L$ is unipotent with fixed line $F\cdot1$, so $e=1$.
\item If $r\equiv1\pmod3$, then $\varphi|_L$ has the two eigenvalues $\omega,\omega^2$ ($\omega\in F$ a primitive cube root of unity), so $e=2$.
\end{itemize}
For a fixed point:
\begin{itemize}
\item If $\varphi x=\omega x$, then $K$ splits into $\varphi$-eigenlines, which $\tau$ preserves, so $\tau x=cx$ and $\tau^{-1}x=c'x$ with $c,c'\in F$, and $\{u:\mathrm{Tr}(xu)=0\}=F+Fx$.
For $u=\beta_0+\beta x$ the condition becomes $3cc'\beta^{r+r^2}x^3=0$, i.e.\ $\beta=0$, i.e.\ $y=x$.
\item If $p=3$, then $L=F+Fv$ with $\varphi v=v+\alpha$, $\alpha\in F^\times$. For $y=\beta_0+\beta v$, the condition reduces to $\mathrm{Tr}(v^2)\,\beta^{r+r^2}=0$ with $\mathrm{Tr}(v^2)=-\alpha^2\ne0$, so again $y=x$.
\end{itemize}
Summing over $x$ gives $N(\rho_1)$ when $m$ is prime, because then every nonzero difference is a unit.
The maps $(a,b,c)\mapsto(b,a,c)$ and $(a,b,c)\mapsto(a,c,b)$ realise $\rho\mapsto1-\rho$ and $\rho\mapsto1/\rho$ and preserve the unit conditions, so $N$ is constant on the $S_3$-orbit.
\end{proof}
The theorem was also checked by independent brute force for $r\le13$; for $q=8,27,512,1331$ it was checked on all six ratios. An independent AI review accepted it after the fixes included here.

\emph{Algebraic form.} The six ratios of the orbit are roots modulo $m$ of $x^3-3x+1$ and $x^3-3x^2+1$. For prime $m$ they are all of the roots;
for composite $m$ these cubics can have more roots, by the Chinese remainder theorem. The first is the minimal polynomial of $2\cos(2\pi/9)=\zeta_9+\zeta_9^{-1}$, with $\zeta_9\mapsto r$.
Its root set is invariant under $x\mapsto1/(1-x)$, the inverse of the order-3 map $s\mapsto1-1/s$ of Theorem~\ref{thm:curve}(a).

\section{The resonance and the cube excess}\label{sec:resblind}
For the six resonant ratios $M(\rho)\approx r\mu$. If all other ratios behaved as in Conjecture~\ref{conj:var}, then $\mathrm{Var}\,M/\mathbb E\,M$ would exceed 1 by about
\[\Delta=6(r-1)^2\mu/n_{\rm good},\qquad n_{\rm good}=\prod_{\ell^e\|m}\ell^{e-1}(\ell-2),\quad \mu=\frac{(q+1)q(q-1)}{3k\,n_{\rm good}}\prod_{\ell\mid m}\frac{(\ell-1)(\ell-2)}{\ell^2}.\]
This is a heuristic, not a theorem; it decays like $q^{-1/3}$.

\emph{Registration.}
\begin{itemize}
\item The value of $\Delta$ was registered before computation only for $q=12167$.
\item The statements "the six largest $M(\rho)$ are roots of the cubics" and "$M/\mu\in[0.8r,1.2r]$" were registered before computation for $q=4913,6859$. For $q=15625$ the registration was "the cubic orbit is among the twelve largest, with $M/\mu\in[20,30]$".
\item All other entries of Table~\ref{tab:res} are retrospective.
\end{itemize}
The plan file and its SHA-256 ledger are released. Registration times are self-attested.
\begin{table}[h]\centering\small
\begin{tabular}{rrrrrl}\toprule $q$&$k$&$\Delta$ formula&observed&top-6 $M/\mu$ ($r$)&status\\ \midrule
125&3&0.089&0.101&4.91 (5)&retrospective\\
343&3&0.143&0.180&7.19 (7)&retrospective\\
512&9&0.021&0.039&8.03 (8)&retrospective\\
729&6&0.032&0.042&8.86 (9)&retrospective\\
1331&3&0.0500&0.0534&11.0 (11)&retrospective (local computation)\\
2197&3&0.0873&0.0877&-- (13)&retrospective\\
4096&12&0.020&0.022&15.93 (16)&retrospective\\
4913&3&0.0368&0.0377&17.01 (17)&top-6, $M/\mu$ registered\\
6859&3&0.0631&0.0635&18.94 (19)&top-6, $M/\mu$ registered; $\Delta$ retrospective\\
12167&3&0.0281&0.0281&22.99 (23)&\textbf{$\Delta$ registered (blind)}; top-6, $M/\mu$ retrospective\\
15625&6&0.0261&0.0261&24.95 (25)&top-6, $M/\mu$ registered; $\Delta$ not registered\\
19683&9&0.0077&0.0092&27.01 (27)&exploratory ($p=3$)\\
\bottomrule\end{tabular}
\caption{The resonance and the cube excess. Wherever the six largest $M(\rho)$ were recorded, they are roots of the two cubics and $M/\mu\approx r$.
Non-cube controls ($q=8192,14641,16384$) give ratio$/3k\in[0.9997,1.0024]$.}\label{tab:res}\end{table}

For $k=3$ and $q\ge1331$, the formula gives $94$--$100\%$ of the observed excess. In the other rows of Table~\ref{tab:res} it gives $55$--$100\%$ (lowest at $q=512$).
Two earlier registered readings failed at $q=6859$: a $q^{-1/2}$ decay band, and a $p\equiv1\pmod3$ effect. Registered tests at $q=24389$ and $q=32768$ are pending.
For $k=6,9,12$ a second, weaker orbit appears, with $M/\mu=2.77,\,3.91,\,4.00,\,2.98$ at $q=729,4096,15625,19683$. Its ratios are not roots of integer
polynomials of degree $\le4$ with coefficients bounded by 4, and they are not of the form $p^a\pm p^b\pm p^c$. We have not identified this orbit.
At $q=19683$ it contributes only about $10\%$ of the residual $0.0015$; most of the residual comes from several hundred ratios with $M/\mu\approx2$.
Weaker levels with $M/\mu\approx2$ also occur for $k=3$ ($q=4913,12167$) and for the non-cube control $q=16384$. They are also unexplained.

\section{What remains}\label{sec:rem}
By Proposition~\ref{prop:red}, a negative answer to the $O(1)$ question follows from $R_k\to0$ for each fixed $k$ and each $C$. For $k=3$ this follows
from Conjecture~\ref{conj:var}, or from any bound $o(q^4/\log q)$ on the variance. For $k\ge4$ a different decomposition is needed. In
Fourier terms the missing input is cancellation, on average over the dilation parameter, in correlations of Jacobi sums over
$\mathbb F_{q^3}$, beyond the square-root (Weil--Deligne) bound for single sums. We do not know how to prove it.

\section{Limitations, literature and disclosure}\label{sec:lim}
Theorem~\ref{thm:curve} covers fixed ratios only (prime $q$, or large characteristic). The higher-moment and variance laws are
numerical. Theorem~\ref{thm:res} is proved; two polynomial identities and one factorisation in its proof were verified by exact computer algebra (Sage 10.7; the scripts and their SHA-256 hashes are released) rather than by hand. The excess formula of Section~\ref{sec:resblind} is a heuristic, not a theorem. For Version~3, a search (web, zbMATH Open, three full texts; 20 October 2026) found no prior statement of Theorem~\ref{thm:res}. Proposition~\ref{prop:orb} is elementary and probably known. Proposition~\ref{prop:red} is conditional on moment convergence, which is not proved. A full-text literature review and a
zbMATH Open search (19 October 2026) found the forbidden ratios, the arc and oval facts \cite{JV,Pi} and the dilate strategy
(\cite{Go}, Shakan \cite{Sh}). They did not find the curve count, the two laws, the identity of Section~6 or the reduction of
Section~4. Guy's C9 entry and MathSciNet were not checked. Computations and drafting were done with AI assistance (Anthropic Claude).
The statements labelled as proved were checked by an independent AI review and by exact computation, but not yet by a human expert.

\section*{Data and code}
Version 1 release: \url{https://github.com/vichienFF/ramanujan-explorer/tree/main/erdos_sidon}. Version 2 additions (certificates,
blind-test plans with SHA-256 ledger, Kaggle outputs): \url{https://github.com/vichienFF/ramanujan-explorer/tree/main/erdos_sidon/v2}. Version 3 additions (Lean files, Sage scripts for Theorem~\ref{thm:res}, resonance blind tests): \url{https://github.com/vichienFF/ramanujan-explorer/tree/main/erdos_sidon/v3}.

\begin{thebibliography}{99}
\bibitem{BHP} R.~C.~Baker, G.~Harman, J.~Pintz, The difference between consecutive primes II, Proc. London Math. Soc. (3) 83 (2001) 532--562.
\bibitem{BFR} J.~Balogh, Z.~F\"uredi, S.~Roy, An upper bound on the size of Sidon sets, Amer. Math. Monthly 130(5) (2023) 437--445.
\bibitem{Bo} R.~C.~Bose, An affine analogue of Singer's theorem, J. Indian Math. Soc. (N.S.) 6 (1942) 1--15.
\bibitem{ET} P.~Erd\H{o}s, P.~Tur\'an, On a problem of Sidon in additive number theory, J. London Math. Soc. 16 (1941) 212--215.
\bibitem{EP} T.~Bloom, Erd\H{o}s problems, problem \#30, \url{https://www.erdosproblems.com/30}.
\bibitem{Go} W.~T.~Gowers, What are dense Sidon subsets of $\{1,2,\dots,n\}$ like?, blog post and comments (2012),
\url{https://gowers.wordpress.com/2012/07/13/what-are-dense-sidon-subsets-of-12-n-like/}.
\bibitem{Guy} R.~K.~Guy, Unsolved Problems in Number Theory, 3rd ed., Springer 2004, problem C9.
\bibitem{MO} G.~Martin, K.~O'Bryant, Constructions of generalized Sidon sets, J. Combin. Theory Ser. A 113 (2006) 591--607.
\bibitem{Sh} G.~Shakan, A large gap in a dilate of a set, arXiv:2004.14828 (2020).
\bibitem{RD} T.~Rokicki, G.~Dogon, Golomb ruler and Sidon set tables, \url{https://cube20.org/golomb/}.
\bibitem{AP} Y.~Aubry, M.~Perret, A Weil theorem for singular curves, in: R.~Pellikaan, M.~Perret, S.~G.~Vl\u{a}du\c{t} (eds.), Arithmetic, Geometry and Coding Theory (Luminy 1993), de Gruyter, Berlin 1996, 1--7.
\bibitem{JV} D.~Jungnickel, K.~Vedder, On the geometry of planar difference sets, European J. Combin. 5 (1984) 143--148.
\bibitem{Pi} G.~Pickert, Differenzmengen und Ovale, Discrete Math. 73 (1989) 165--179.
\bibitem{Gr} B.~Green, 100 open problems, Problems 31--32 (manuscript, updated 2026).
\bibitem{Ha} M.~Hall, Jr., Cyclic projective planes, Duke Math. J. 14 (1947) 1079--1090.
\bibitem{MR} R.~L.~McFarland, B.~F.~Rice, Translates and multipliers of abelian difference sets, Proc. Amer. Math. Soc. 68 (1978) 375--379.
\bibitem{Si} J.~Singer, A theorem in finite projective geometry and some applications to number theory, Trans. Amer. Math. Soc. 43 (1938) 377--385.
\bibitem{Z} Z.~Zhang, Finding finite $B_2$-sequences with larger $m-a_m^{1/2}$, Math. Comp. 63 (1994) 403--414.
\end{thebibliography}
\end{document}
